\section{Introduction}\label{sec:introduction}

Let $E$ be an elliptic curve over $\Q$. Write
\[
 a(E)=\ord_{s=1}L(E,s),\qquad r(E)=\rank_{\Z}E(\Q),\qquad
 c_2(E)=\corank_{\Ztwo}\Sel_{2^\infty}(E/\Q).
\]
The Selmer group uses the local Kummer conditions at every place.
Its corank includes both the Mordell--Weil rank and any divisible
$2$-primary contribution from the Tate--Shafarevich group.
Modularity supplies the analytic continuation and functional equation
of $L(E,s)$ \citep{BreuilConradDiamondTaylor2001}.

\begin{theorem}[The $2$-converse]\label{thm:main-converse}
For every elliptic curve $E/\Q$ with $c_2(E)\in\{0,1\}$,
\[
 a(E)=r(E)=c_2(E),\qquad \Sha(E/\Q)\text{ is finite}.
\]
\end{theorem}

There is no restriction on reduction type, complex multiplication,
rational torsion, or rational isogenies. The hypothesis concerns the full
$2$-power Selmer corank; a Mordell--Weil rank bound alone does not give it.

For a nonzero squarefree integer $d$, let $E^{(d)}$ denote the quadratic
twist of $E$, and set
\[
 \mathcal D(X)=\{d\in\Z:0<|d|\le X,\ d\text{ squarefree}\}.
\]
Thus positive and negative parameters are counted together and ordered
by absolute value. A residue class can fix the sign of the functional
equation; the half-density statement below concerns the full set
$\mathcal D(X)$.

\begin{theorem}[Goldfeld's analytic density conjecture]\label{thm:main-density}
For every elliptic curve $E/\Q$ and each $j\in\{0,1\}$,
\[
 \lim_{X\to\infty}
 \frac{\#\{d\in\mathcal D(X):a(E^{(d)})=j\}}{\#\mathcal D(X)}
 =\frac12.
\]
Consequently the twists of analytic rank at least two have density zero.
For a density-one set of signed squarefree parameters $d$, the analytic
and Mordell--Weil ranks of $E^{(d)}$ agree and its whole
Tate--Shafarevich group is finite.
\end{theorem}

The theorem follows by combining the pointwise $2$-converse with
Smith's distribution of full $2$-power Selmer coranks
\citep[Theorem~1.1]{Smith2025}. Section~\ref{sec:conclusions} gives
the deduction, including the passage from integer to squarefree
parameters. The main work of this paper is the pointwise converse.

\subsection{Consequences and companion results}

The converse has a finite-descent consequence. Put
\[
 d_2(E)=\dim_{\F_2}\Sel_2(E/\Q)-\dim_{\F_2}E(\Q)[2].
\]

\begin{corollary}[A finite Selmer criterion]\label{cor:finite-descent}
If $d_2(E)\in\{0,1\}$, then
\[
 a(E)=r(E)=c_2(E)=d_2(E),\qquad
 \Sha(E/\Q)\text{ is finite},\qquad \Sha(E/\Q)[2^\infty]=0.
\]
\end{corollary}

Thus a finite $2$-descent can certify the analytic rank and the vanishing
of the entire $2$-primary Tate--Shafarevich group in these two cases.
The proof uses the alternating Cassels--Tate pairing to pass from
$d_2(E)$ to the full Selmer corank, and is given in
Section~\ref{sec:conclusions}.

The density theorem resolves the frequency assertion for analytic ranks
zero and one. The mean analytic-rank assertion requires additional
control: a set of density zero can still contribute to the rank-weighted
average. The separate companion
\citet[Theorems~1.1 and~1.2]{OpenAIMeanAnalyticRank2026} proves the
required analytic tail estimate and, together with
Theorem~\ref{thm:main-density}, obtains mean analytic rank $1/2$
in the same signed squarefree ordering.

A subsequent companion \citep[Theorem~1.1]{OpenAIAllPrimeSelmer2026}
extends the pointwise converse to every prime $p$: full $p$-power Selmer
corank $r\in\{0,1\}$ forces analytic and Mordell--Weil ranks $r$
and finiteness of the whole Tate--Shafarevich group. Its proof uses the
$2$-converse and analytic twist densities established here.

The rank and finiteness assertions are distinct from the
Birch--Swinnerton-Dyer leading-coefficient formula.
The pointwise two-primary formula of
\citet[Theorem~1.1]{OpenAITwoPrimaryBSD2026}, combined with the
Selmer-corank densities in the proof of Theorem~\ref{thm:main-density},
gives the exact two-primary formula on a density-one set of signed
squarefree twists. Their common analytic and Mordell--Weil ranks are
zero and one with density $1/2$ each. This consequence is
\citet[Corollary~16.1]{OpenAITwoPrimaryBSD2026}; its leading-coefficient
assertion concerns the valuation at $2$.

\subsection{Goldfeld's conjecture and analytic nonvanishing}

Birch and Swinnerton-Dyer's comparisons of rational points with central
$L$-values led to the conjectural equality of algebraic and analytic rank
\citep{BirchSwinnertonDyer1965}. The conjecture also predicts finiteness
of $\Sha(E/\Q)$ and a formula for the leading coefficient
\citep[\S1, pp.~415--421]{Tate1966}.
Goldfeld's original 1979 formulation asks for the mean analytic order
in a family parametrized by quadratic-field discriminants
\citep[p.~113, Conjecture~(B)]{Goldfeld1979}.
The density formulation predicts that analytic ranks zero and one each
occur half the time. Our counting convention throughout is the signed
squarefree set $\mathcal D(X)$.

The analytic study of quadratic twists produced both nonzero central
values and simple central zeros. Bump--Friedberg--Hoffstein and
Murty--Murty supplied auxiliary twists with the nonvanishing needed in
Heegner-point arguments; Iwaniec and Perelli--Pomyka\l a developed
quantitative estimates for central derivatives
\citep{BFH1990,MurtyMurty1991,Iwaniec1990,PerelliPomykala1997}.
Ono--Skinner obtained a lower bound of order $X/\log X$ for the
number of nonvanishing central values among twists by fundamental
discriminants $D$ with $0<|D|\le X$
\citep[Corollary~3]{OnoSkinner1998}.
More recently, Kumar--Mallesham--Sharma--Singh sharpened derivative
nonvanishing estimates through second moments, and
Kriz--Nordentoft improved central-value lower bounds for curves with
no rational two-torsion
\citep[Theorem~1]{KumarEtAl2024}
\citep[Theorem~1.1 and Remark~1.2]{KrizNordentoft2025}.
These quantitative results yield many twists without by themselves
giving a positive proportion for every fixed curve.

Positive proportions were obtained in special families. James proved
analytic-rank-zero proportions for particular curves, and Vatsal proved
positive proportions of both analytic ranks zero and one for $X_0(19)$
\citep[Theorem~1]{James1998} \citep[Theorem~II]{Vatsal1998}.
Kriz--Li proved both positive proportions for every rational elliptic
curve admitting a rational $3$-isogeny
\citep[Theorem~1.5]{KrizLi2019}.
In the congruent-number family, Tian's Heegner-point work gives rank
one and odd-order Tate--Shafarevich group in specified subfamilies
\citep{Tian2014}. Genus-period and genus-point congruences, followed by
Selmer statistics, give further analytic density results
\citep{TianYuanZhang2017,SmithCongruent2016}.
These arguments connect coefficient congruences and arithmetic
statistics with analytic nonvanishing, a connection also central to
the present proof.

The statistical input used here has its own progression.
Smith first established higher $2$-power Selmer distributions for curves
with full rational two-torsion and no rational cyclic subgroup of order
four \citep{Smith2017}. His subsequent work treated the irreducible
residual case and further rational two-isogeny cases
\citep{SmithDistributionI,SmithDistributionII}.
The unrestricted theorem gives full Selmer coranks zero and one with
density $1/2$ each for every rational elliptic curve
\citep[Theorem~1.1]{Smith2025}. Its analytic corollary assumes
Birch--Swinnerton-Dyer \citep[Corollary~1.2]{Smith2025}.
We use the Selmer theorem itself and supply the pointwise analytic
implication through Theorem~\ref{thm:main-converse}.

\subsection{The Selmer converse problem}

Gross--Zagier and Kolyvagin established the forward implication:
analytic rank at most one gives equality with the Mordell--Weil rank
and finiteness of the whole Tate--Shafarevich group
\citep{GrossZagier1986,Kolyvagin1990}; see also the combined formulation
in \citet[\S1]{Skinner2020}.
The Gross--Zagier formula relates a central derivative to the height
of a Heegner point, and Kolyvagin's derived classes control the Selmer
group. A converse must produce analytic nonvanishing from Selmer
information.

Cyclotomic Iwasawa theory gave rank-zero converses at good ordinary
odd primes in the work of Skinner--Urban, under residual
irreducibility and ramification hypotheses
\citep{SkinnerUrban2014}.
Skinner and Wei Zhang proved rank-one converse criteria through
Heegner points and derived classes, under their respective reduction,
residual-image and local hypotheses \citep{Skinner2020,Zhang2014}.
Burungale--Castella--Skinner subsequently removed the auxiliary
ramification conditions in the good ordinary, residually irreducible
case for $p>3$ \citep[Theorem~1.2.2 and the discussion following
Remark~1.2.3]{BurungaleCastellaSkinner2025}.
At good supersingular primes $p>3$, Castella--Wan obtained the rank-one
Selmer-corank implication for semistable curves
\citep[Theorem~A]{CastellaWan2024}.

Reducible residual representations call for a different Iwasawa-theoretic
comparison. Castella--Grossi--Lee--Skinner proved the corank-zero and
corank-one converses at good ordinary Eisenstein primes $p>2$, with
local isogeny character neither trivial nor cyclotomic
\citep[Theorem~E]{CastellaGrossiLeeSkinner2022}.
Keller--Yin treated potentially good ordinary reduction at odd
Eisenstein primes without that character exclusion
\citep[Theorem~B]{KellerYin2024}.
These results enlarge the allowed residual representations while
retaining conditions on the prime and its reduction type.

The CM theory includes substantial results at the prime two.
Burungale--Tian proved the rank-one converse at good ordinary primes
$p>3$ \citep{BurungaleTian2020};
Burungale--Castella--Skinner--Tian included the small ordinary primes
under an additional conductor hypothesis
\citep[Theorem~A]{BurungaleCastellaSkinnerTian2022}.
Ressler and Yu obtained the small-prime good-ordinary result without
that conductor condition; see Yu's Theorem~1 and the attribution in
Burungale--Castella--Skinner--Tian, Remark~C
\citep{Yu2021,BurungaleCastellaSkinnerTian2022}.
Burungale--Skinner's combined ordinary formulation includes $p=2$,
and their good supersingular CM theorem also permits $p=2$ under an
additional surjectivity condition on Selmer localization at $p$
\citep[Theorems~A.3 and~A.1]{BurungaleSkinner2024}.

In corank zero, Burungale--Tian's theorem covers every rational CM
elliptic curve at every prime, without a reduction restriction
\citep[Theorem~1.1]{BurungaleTian2026}.
Combined with Smith's unrestricted distribution, it therefore already
gives the analytic-rank-zero half for every CM twist family.
Their Theorem~1.2 gives rank zero for almost all positive squarefree
$n\equiv1,2,3\pmod8$ in the congruent-number family
$ny^2=x^3-x$.
Kriz's preprint states a converse for curves with CM by the full ring
of integers at the unique ramified prime of the CM field, and derives
analytic densities for restricted CM families, including congruent-number
curves \citep[Theorems~1.1 and~10.17]{Kriz2022}.
Its dyadic pointwise statement concerns the CM fields $\Q(i)$ and
$\Q(\sqrt{-2})$.
Together with Smith's distribution theorem, the stated dyadic converse
would yield both half-density conclusions for quadratic twists of every
curve with CM by either of these maximal orders.
Theorem~\ref{thm:main-converse} treats the full $2$-power Selmer corank
for every rational elliptic curve, in both ranks zero and one.

\subsection{Proof strategy}

Fix the curve under consideration. We construct finite binary families
of twists
\[
 h_x=\prod_{q\in\mathcal Q}(q^*)^{\lambda_q(x)},\qquad
 q^*=(-1/q)q,\qquad x\in\F_2^b,
\]
where $\mathcal Q$ is a finite set of distinct good odd primes and the
$\lambda_q$ are linear forms. The parameter at $x=0$ is $1$.
The families have the same local twisting classes at every fixed bad
place. Their nonzero vertices have analytic order zero or one, according
to the corank of the fixed curve. A scalar determinant coordinate
attached to the relevant class and Selmer complex has a nonzero central
specialization there, with uniformly bounded $2$-adic valuation.
Bounded integral interpolation then recovers nonvanishing at the missing
vertex $0$.

At corank zero, the interpolated class is a Beilinson--Kato class and
its central specialization detects $L(E,1)$ \citep{Kato2004}.
At corank one, it is a Heegner class over an imaginary quadratic field
$K=\Q(\sqrt{k})$,
whose height detects
\[
 \left[L(E,s)L(E^{(k)},s)\right]'_{s=1}.
\]
An analytically nonvanishing companion $E^{(k)}$ makes the Selmer
corank over $K$ equal to one. The argument retains integral Kummer
conditions at $2$ and every bad prime. The rational Iwasawa input is
used in its stated scope; the uniform integral estimates at $2$ are
proved in Sections~\ref{sec:heegner-interpolation} and
\ref{sec:cy:interpolation}. The Heegner construction is related to
Howard's Kolyvagin-system framework \citep{Howard2012}; the integral
bounds required here are established with the local conditions and
uniformity stated in this paper.

The finite families are built from modular coefficients. Waldspurger's
formula supplies the even-sign coefficients \citep{Waldspurger1981}.
In odd sign, reductions
of genus Heegner sums give weighted ternary theta coefficients at
increasing finite precisions. A uniform lower bound allows a least
normalized valuation to be selected. In the even case this is a minimum
among the coefficients of one form. In the odd case it is a minimum
over bounded-weight sequences as the precision increases; each finite
collection of tests is then realized at a sufficiently high actual
precision. A test attaining the minimum detects the required analytic
order and gives a uniform upper bound. Weight-two Hecke traces turn these coefficient
tests into contraction and deletion rules for finite configurations
of primes.

There are two exhaustive residual cases. If $E[2]$ is reducible, its
invariant line is a rational point of order two. Graphs encode prescribed
quadratic residue symbols between auxiliary primes. The resulting finite
construction makes two separate coefficient tests nonzero together at
every nonzero binary vertex. The odd argument fixes one companion field
before the dimension grows. In this branch, the even construction first
supplies the lower bound needed to normalize the odd coefficient test.
If $E[2]$ is irreducible, local Selmer equations give matrices over
$\F_2$. Subtracting a correction determined by finitely many local types
makes these matrices symmetric; their nullities still give a lower
bound on the Selmer dimension up to a fixed cost. Removing one selected prime or pair of primes, or retaining
it with either of two allowed local choices, gives three configurations.
Their coefficient tests satisfy the same sum-zero relation as the
corresponding matrix determinants. This relation does not identify the
coefficient test with the determinant.
We group the configurations into blocks and arrange that the matrix
entries between distinct blocks vanish.
A nonzero coefficient test bounds the matrix kernel, and hence the number
of singular blocks in the finite configuration. Fixing a configuration
with a maximal number of such blocks lets us arrange nonzero tests at
every nonzero binary vertex. An additional alternating-matrix construction
supplies a companion with uniformly bounded prime count and Selmer length.

Two parts of the argument have uses beyond this assembly. The
interpolation algebra clears denominators through a bounded complex,
so the required congruence precision is independent of the number of
new primes. The graph construction treats the coefficient tests as polynomials in
prescribed residue symbols and makes separately nonzero tests nonzero
together. Its proof retains squared variables until after multiplying
their highest-degree parts.
Neither construction presupposes the desired nonvanishing at the base.

\paragraph{Organization.}
Section~\ref{sec:preliminaries} records the rank implications and local
conventions. Section~\ref{sec:transfer-criteria} states the two arithmetic
transfer criteria, the lower bounds needed to normalize coefficients,
and auxiliary nonvanishing with prescribed local conditions.
Sections~\ref{cs:coefficient-systems} and~\ref{sec:trace-calculus}
turn analytic nonvanishing into coefficient tests and prime-substitution
identities. Section~\ref{gr:section} develops the graph construction
used in the rational-torsion converse of Section~\ref{pw:section}.
Section~\ref{sec:lm:local-matrices} supplies the Selmer matrices for
the irreducible converse in Section~\ref{sec:sa:saturation}.
The remaining main-text sections prove the arithmetic inputs:
Section~\ref{sec:auxiliary-nonvanishing} establishes auxiliary
nonvanishing, Section~\ref{sec:binary-algebra} develops the common
finite-complex and binary transfer arguments, and
Sections~\ref{sec:cy:interpolation} and~\ref{sec:heegner-interpolation}
construct the cyclotomic and Heegner determinant coordinates and prove
their integral bounds. Finally, Section~\ref{sec:conclusions} assembles the unrestricted
result and its density and finite-descent consequences.
